\documentclass[a4paper,12pt]{article}

\usepackage[T2A]{fontenc}
\usepackage[utf8]{inputenc}
\usepackage[russian]{babel}
\usepackage{amsmath,amssymb}
\usepackage{paratype}
\renewcommand{\familydefault}{\sfdefault}
\usepackage{icomma}
\usepackage[a4paper,left=20mm,right=20mm,top=20mm,bottom=22mm]{geometry}
\usepackage{microtype}
\usepackage{tikz}
\usetikzlibrary{calc,arrows.meta,positioning,shapes.geometric}
\usepackage{tkz-euclide}
\usepackage{tabularx,booktabs}
\usepackage{enumitem}
\usepackage{parskip}
\usepackage{fancyhdr}
\usepackage{setspace}
\usepackage{xcolor}

\definecolor{accent}{HTML}{008080}
\definecolor{darkaccent}{HTML}{004D4D}
\definecolor{textgray}{HTML}{303434}
\definecolor{linegray}{HTML}{8A9696}

\color{textgray}
\onehalfspacing
\setlength{\parindent}{0pt}
\setlength{\emergencystretch}{2em}

\setlist[itemize]{leftmargin=1.55em,itemsep=0.28em,topsep=0.35em}
\setlist[enumerate]{leftmargin=1.7em,itemsep=0.38em,topsep=0.4em}

\pagestyle{fancy}
\fancyhf{}
\fancyhead[L]{\small Григорий Архипов · 10 класс}
\fancyhead[R]{\small Математика · чек-лист}
\fancyfoot[C]{\small\thepage}
\renewcommand{\headrulewidth}{0.35pt}
\renewcommand{\footrulewidth}{0pt}
\setlength{\headheight}{14pt}

\newcommand{\sectionline}{%
  \vspace{-0.35em}\par\noindent
  \textcolor{accent!55}{\rule{\linewidth}{0.45pt}}\vspace{0.35em}
}

\newcommand{\keyidea}[2]{%
  \par\smallskip
  \noindent\textcolor{darkaccent}{\textbf{#1}} #2
  \par\smallskip
}

\newcommand{\checkitem}{\item[$\square$]}

\tikzset{
  flowstart/.style={
    ellipse,
    draw=accent!75,
    fill=accent!5,
    minimum width=4.4cm,
    minimum height=9mm,
    align=center,
    font=\small
  },
  flowaction/.style={
    rectangle,
    rounded corners=1.5mm,
    draw=accent!65,
    fill=white,
    text width=7.0cm,
    minimum height=9mm,
    align=center,
    font=\small
  },
  flowdecision/.style={
    diamond,
    aspect=2.2,
    draw=darkaccent!80,
    fill=accent!4,
    text width=4.2cm,
    inner sep=1.3pt,
    align=center,
    font=\small
  },
  flowend/.style={
    ellipse,
    draw=darkaccent!80,
    fill=accent!7,
    minimum width=4.4cm,
    minimum height=9mm,
    align=center,
    font=\small
  },
  flowarrow/.style={
    -{Stealth[length=3mm,width=2mm]},
    line width=0.8pt,
    draw=textgray
  },
  flowlabel/.style={
    font=\small,
    fill=white,
    inner sep=1.5pt
  }
}

\begin{document}

% ============================================================
% ТИТУЛЬНЫЙ ЛИСТ
% ============================================================

\begin{titlepage}
\thispagestyle{empty}
\centering

\vspace*{24mm}

{\Large\bfseries\textcolor{darkaccent}{ЧЕК-ЛИСТ}\par}

\vspace{8mm}

{\huge\bfseries Понижение степени многочлена\par}

\vspace{3mm}

{\Large\bfseries Теорема Безу и схема Горнера\par}

\vspace{8mm}

{\Large\bfseries Построение сечений многогранников\par}

\vspace{18mm}

{\large Математика · 10 класс\par}

\vspace{4mm}

{\large 17 сентября 2026 года\par}

\vfill

{\large Лёвин Артём Александрович\par}

\vspace{2mm}

{\normalsize эксклюзивно для Григория Архипова\par}

\vspace{21mm}

\begin{tikzpicture}[remember picture,overlay]
  \draw[
    accent!45,
    line width=1.05pt
  ]
    ([xshift=22mm,yshift=18mm]current page.south west)
    .. controls
    ([xshift=66mm,yshift=34mm]current page.south west)
    and
    ([xshift=-72mm,yshift=9mm]current page.south east)
    ..
    ([xshift=-22mm,yshift=22mm]current page.south east);
\end{tikzpicture}

\end{titlepage}

% ============================================================
% ОСНОВНОЙ ЧЕК-ЛИСТ
% ============================================================

\section*{1. Что важно уметь после занятия}
\sectionline

На занятии разбирались две самостоятельные темы:
понижение степени многочлена и построение сечений многогранников.
В первой теме основной инструмент --- схема Горнера после нахождения корня.
Во второй --- поиск линий пересечения секущей плоскости с гранями
с помощью общих точек, следов и проекций.

\begin{itemize}
  \checkitem
  Я умею находить возможные рациональные корни многочлена
  и проверять их подстановкой.

  \checkitem
  Я понимаю связь $P(r)=0$ и делимости $P(x)$ на $(x-r)$.

  \checkitem
  Я умею понижать степень делением уголком и схемой Горнера.

  \checkitem
  Я записываю нулевой коэффициент,
  если какая-либо степень $x$ отсутствует.

  \checkitem
  Я различаю грань многогранника и плоскость этой грани.

  \checkitem
  Я соединяю две точки только после установления общей плоскости.

  \checkitem
  Я умею получать вспомогательные точки
  на продолжениях рёбер и прямых.

  \checkitem
  Я понимаю,
  как проекции помогают найти след секущей плоскости на основании.
\end{itemize}

\keyidea{Общая стратегия.}{
Сначала ищите самый короткий корректный путь.
Для многочлена проверьте вынесение общего множителя,
группировку и знакомые формулы.
Для сечения сначала определите,
какие известные точки уже принадлежат одной грани.
}

% ============================================================
% АЛГЕБРА
% ============================================================

\section*{2. Понижение степени многочлена}
\sectionline

Пусть
\[
P(x)
=
a_nx^n
+
a_{n-1}x^{n-1}
+
\dots
+
a_1x
+
a_0,
\qquad
a_n\ne0.
\]

Число $r$ называется корнем многочлена $P(x)$, если
\[
P(r)=0.
\]

По теореме Безу остаток от деления $P(x)$ на $(x-r)$ равен $P(r)$.
Поэтому
\[
P(r)=0
\quad\Longleftrightarrow\quad
P(x)=(x-r)Q(x),
\]

\keyidea{Зачем это нужно.}{
Найденный корень даёт линейный множитель $(x-r)$
и позволяет уменьшить степень уравнения на единицу.
}

\subsection*{2.1. Где искать рациональный корень}

Пусть все коэффициенты $P(x)$ целые и
\[
r=\frac pq,
\qquad
(p,q)=1,
\]
--- рациональный корень.
Тогда
\[
p\mid a_0,
\qquad
q\mid a_n.
\]

Для приведённого многочлена $a_n=1$,
поэтому любой рациональный корень
является целым делителем свободного коэффициента $a_0$.

Например, для
\[
x^3-6x^2+11x-6=0
\]
достаточно проверить
\[
\pm1,\quad
\pm2,\quad
\pm3,\quad
\pm6.
\]

При $x=1$:
\[
1-6+11-6=0,
\]
следовательно, $x=1$ --- корень,
а $(x-1)$ --- множитель.

\subsection*{2.2. Деление уголком: смысл понижения степени}

Разделим
\[
x^3-6x^2+11x-6
\]
на $(x-1)$.

Последовательно получаем
\[
\frac{x^3}{x}=x^2,
\qquad
(x^3-6x^2)-(x^3-x^2)=-5x^2,
\]
\[
\frac{-5x^2}{x}=-5x,
\qquad
(-5x^2+11x)-(-5x^2+5x)=6x,
\]
\[
\frac{6x}{x}=6,
\qquad
(6x-6)-(6x-6)=0.
\]

Итак,
\[
x^3-6x^2+11x-6
=
(x-1)(x^2-5x+6).
\]

После разложения квадратного трёхчлена:
\[
x^3-6x^2+11x-6
=
(x-1)(x-2)(x-3),
\]
поэтому
\[
\boxed{x=1,\;2,\;3}.
\]

\subsection*{2.3. Схема Горнера}

Схема Горнера выполняет то же деление на $(x-r)$
через коэффициенты многочлена.

Для
\[
x^3-5x^2+2x+8=0
\]
проверяем $r=-1$:
\[
P(-1)=-1-5-2+8=0.
\]

Записываем коэффициенты
\[
1,\;-5,\;2,\;8
\]
и выполняем последовательность
«умножить на $r$ --- сложить»:
\[
\begin{array}{c@{\quad}rrrr}
-1 & 1 & -5 & 2 & 8\\
   &   & -1 & 6 & -8\\
\cmidrule(lr){2-5}
   & 1 & -6 & 8 & 0
\end{array}
\]

Первые три числа нижней строки --- коэффициенты частного:
\[
Q(x)=x^2-6x+8,
\]
последнее число --- остаток.

Поэтому
\[
x^3-5x^2+2x+8
=
(x+1)(x^2-6x+8)
=
(x+1)(x-2)(x-4),
\]
откуда
\[
\boxed{x=-1,\;2,\;4}.
\]

\keyidea{Самопроверка.}{
Если $r$ действительно является корнем,
последнее число в схеме Горнера равно нулю.
Ненулевой остаток требует повторной проверки
подстановки и арифметики в таблице.
}

\subsection*{2.4. Пропущенная степень}

В таблице Горнера нужны коэффициенты всех степеней
от старшей до нулевой.
Пропущенной степени соответствует коэффициент $0$.

Например,
\[
x^4-1
=
x^4+0x^3+0x^2+0x-1,
\]
поэтому строка коэффициентов имеет вид
\[
1,\;0,\;0,\;0,\;-1.
\]

Пропуск любого из нулей
сдвигает позиции коэффициентов
и разрушает вычисление.

\subsection*{2.5. Рабочий порядок}

\begin{enumerate}
  \item
  Запишите многочлен по убыванию степеней
  и восстановите нулевые коэффициенты.

  \item
  Проверьте простые способы разложения:
  общий множитель,
  группировку,
  формулы,
  удобную замену.

  \item
  Если требуется подбор корня,
  составьте список допустимых кандидатов.

  \item
  Начинайте проверку с небольших по модулю чисел:
  обычно $\pm1$,
  затем $\pm2$ и т.~д.

  \item
  После найденного корня
  примените схему Горнера или деление уголком.

  \item
  Решите полученное уравнение меньшей степени;
  при необходимости повторите процедуру.
\end{enumerate}

% ============================================================
% БЛОК-СХЕМА АЛГЕБРЫ
% ============================================================

\clearpage
\thispagestyle{empty}

\begin{center}
{\Large\bfseries\textcolor{darkaccent}{
Алгоритм: рациональный корень и схема Горнера
}}
\end{center}

\vspace{3mm}

\begin{center}
\begin{tikzpicture}[
  node distance=3.2mm and 13mm
]

\node[flowstart]
(start)
{Дано уравнение $P(x)=0$};

\node[
  flowaction,
  below=of start
]
(standard)
{
Запишите коэффициенты по убыванию степеней;
пропуски замените нулями
};

\node[
  flowaction,
  below=of standard
]
(candidates)
{
Составьте список допустимых рациональных кандидатов
$r=\frac pq$
};

\node[
  flowaction,
  below=of candidates
]
(test)
{
Проверьте очередной кандидат:
вычислите $P(r)$
};

\node[
  flowdecision,
  below=5mm of test
]
(root)
{$P(r)=0$?};

\node[
  flowaction,
  below=5mm of root
]
(horner)
{
Выполните схему Горнера
для найденного $r$
};

\node[
  flowdecision,
  below=5mm of horner
]
(rest)
{Остаток равен $0$?};

\node[
  flowaction,
  below=5mm of rest
]
(lower)
{
Получите частное $Q(x)$
степени на единицу меньше
};

\node[
  flowdecision,
  below=5mm of lower
]
(easy)
{$Q(x)=0$ уже решается известным способом?};

\node[
  flowaction,
  below=5mm of easy
]
(solve)
{
Решите оставшееся уравнение
и соберите все корни
};

\node[
  flowend,
  below=of solve
]
(end)
{Проверьте корни подстановкой};

\draw[flowarrow]
(start) -- (standard);

\draw[flowarrow]
(standard) -- (candidates);

\draw[flowarrow]
(candidates) -- (test);

\draw[flowarrow]
(test) -- (root);

\draw[flowarrow]
(root)
--
node[flowlabel,right]{да}
(horner);

\draw[flowarrow]
(root.west)
--
++(-1.0,0)
|-
node[
  flowlabel,
  pos=0.18,
  left
]
{нет: следующий $r$}
(test.west);

\draw[flowarrow]
(horner) -- (rest);

\draw[flowarrow]
(rest)
--
node[flowlabel,right]{да}
(lower);

\draw[flowarrow]
(rest.east)
--
++(1.0,0)
|-
node[
  flowlabel,
  pos=0.18,
  right
]
{нет: проверить вычисления}
(test.east);

\draw[flowarrow]
(lower) -- (easy);

\draw[flowarrow]
(easy)
--
node[flowlabel,right]{да}
(solve);

\draw[flowarrow]
(easy.west)
--
++(-1.0,0)
|-
node[
  flowlabel,
  pos=0.16,
  left
]
{нет: повторить для $Q$}
(candidates.west);

\draw[flowarrow]
(solve) -- (end);

\end{tikzpicture}
\end{center}

\clearpage

% ============================================================
% ГЕОМЕТРИЯ
% ============================================================

\section*{3. Сечения многогранников}
\sectionline

Секущая плоскость $\alpha$
пересекает грани многогранника по отрезкам.
Эти отрезки образуют многоугольник ---
сечение многогранника плоскостью $\alpha$.

\keyidea{Главное правило.}{
Если две точки $X$ и $Y$
принадлежат и секущей плоскости $\alpha$,
и плоскости одной грани $\pi$,
то прямая $XY$ является линией пересечения
плоскостей $\alpha$ и $\pi$.
Внутри грани нужная часть этой прямой
даёт сторону сечения.
}

\subsection*{3.1. Что разрешает каждое построение}

\begin{itemize}
  \item
  Если две известные точки сечения
  лежат на одной грани,
  их можно соединить.

  \item
  Продолжения рёбер и сторон сечения
  можно использовать для получения
  вспомогательных точек вне многогранника.

  \item
  Пересекать две прямые допустимо
  после установления их принадлежности
  одной плоскости.

  \item
  Ребро многогранника принадлежит двум соседним граням,
  поэтому часто служит переходом
  от одной грани к другой.

  \item
  Для обоснования стороны сечения
  достаточно указать две её точки
  и общую грань,
  в которой они лежат.
\end{itemize}

\subsection*{3.2. След секущей плоскости}

След плоскости $\alpha$
на некоторой плоскости $\pi$ --- прямая
\[
\alpha\cap\pi.
\]

Чтобы построить эту прямую,
достаточно получить две различные точки,
принадлежащие обеим плоскостям.

Особенно полезен след
на плоскости основания.
Его пересечения с прямыми,
содержащими рёбра основания,
дают вспомогательные точки
для перехода на боковые грани.

\subsection*{3.3. Зачем нужны проекции}

Сложный случай возникает,
когда среди трёх заданных точек
нет пары,
лежащей на одной грани.
Тогда удобно спроецировать две точки
на основание в одном направлении.

Пусть $L_0$ и $M_0$ ---
параллельные проекции точек $L$ и $M$
на плоскость основания.

Прямые $LL_0$ и $MM_0$ параллельны,
поэтому точки
\[
L,\ M,\ M_0,\ L_0
\]
лежат в одной плоскости.

Если
\[
T=LM\cap L_0M_0,
\]
то
\[
T\in LM\subset\alpha
\quad\text{и}\quad
T\in L_0M_0\subset\pi_{\text{осн}}.
\]

Следовательно,
$T$ --- точка следа секущей плоскости
на основании.

Если на основании уже есть точка
$K\in\alpha$,
то прямая
\[
KT
\]
и есть след $\alpha$
на плоскости основания.

Если
\[
LM\parallel L_0M_0,
\]
выбирают другую пару точек
или другой вспомогательный ход.

\subsection*{3.4. Пример с кубом: три точки на разных рёбрах}

В кубе
\[
ABCDA_1B_1C_1D_1
\]
возьмём середины
\[
K\in AB,
\qquad
L\in B_1C_1,
\qquad
M\in DD_1.
\]

Никакая пара $K,L,M$
не лежит в одной грани,
поэтому сразу провести сторону сечения нельзя.

Проецируем $L$ и $M$
вертикально на основание:
\[
L_0\in BC,
\qquad
M_0=D.
\]

Строим
\[
T=LM\cap L_0D.
\]

Точка $T$ принадлежит
и секущей плоскости,
и плоскости основания.

Вместе с $K$
она задаёт след секущей плоскости на основании:
\[
t=KT.
\]

Продолжая $t$,
получаем вспомогательные точки
на прямых $BC$ и $CD$.
Через них последовательно переходим
на правую и заднюю грани.

В результате секущая плоскость
проходит через середины шести рёбер
\[
AB,\quad
BB_1,\quad
B_1C_1,\quad
C_1D_1,\quad
DD_1,\quad
AD.
\]

Сечение --- шестиугольник.

\begin{center}
\begin{tikzpicture}[scale=0.95]

\tkzSetUpLine[
  line width=0.85pt,
  color=linegray
]

\tkzSetUpPoint[
  color=darkaccent,
  fill=darkaccent,
  size=3
]

\tkzSetUpLabel[
  color=black,
  font=\small
]

\tkzDefPoint(0,0){A}
\tkzDefPoint(5,0){B}
\tkzDefPoint(7,1.5){C}
\tkzDefPoint(2,1.5){D}

\tkzDefPoint(0,4){A1}
\tkzDefPoint(5,4){B1}
\tkzDefPoint(7,5.5){C1}
\tkzDefPoint(2,5.5){D1}

\tkzDefPoint(2.5,0){K}
\tkzDefPoint(5,2){N}
\tkzDefPoint(6,4.75){L}
\tkzDefPoint(4.5,5.5){Q}
\tkzDefPoint(2,3.5){M}
\tkzDefPoint(1,0.75){P}

\tkzDrawSegments(
  A,B
  B,C
  C,D
  D,A
  A,A1
  B,B1
  C,C1
  D,D1
  A1,B1
  B1,C1
  C1,D1
  D1,A1
)

\draw[
  accent,
  line width=1.75pt
]
(K)--(N)--(L)--(Q)--(M)--(P)--cycle;

\tkzDrawPoints(K,N,L,Q,M,P)

\tkzLabelPoint[below left](A){$A$}
\tkzLabelPoint[below](B){$B$}
\tkzLabelPoint[right](C){$C$}
\tkzLabelPoint[left](D){$D$}
\tkzLabelPoint[left](A1){$A_1$}
\tkzLabelPoint[above](B1){$B_1$}
\tkzLabelPoint[right](C1){$C_1$}
\tkzLabelPoint[above left](D1){$D_1$}
\tkzLabelPoint[below](K){$K$}
\tkzLabelPoint[right](N){$N$}
\tkzLabelPoint[right](L){$L$}
\tkzLabelPoint[above](Q){$Q$}
\tkzLabelPoint[left](M){$M$}
\tkzLabelPoint[left](P){$P$}

\end{tikzpicture}
\end{center}

\keyidea{Как проверять готовое сечение.}{
Идите по его сторонам по кругу
и для каждой пары соседних вершин
называйте общую грань.
Если для какого-либо отрезка
такой грани нет,
построение требует исправления.
}

\subsection*{3.5. След в пирамиде}

В пирамиде $SABCD$ пусть
\[
K\in SA,
\qquad
M\in SB,
\qquad
N\in SC.
\]

Если соответствующие прямые не параллельны,
для нахождения точки сечения
на ребре $SD$ можно построить
\[
F=KM\cap AB,
\qquad
Q=MN\cap BC.
\]

Точки $F$ и $Q$
лежат в плоскости основания
и одновременно принадлежат секущей плоскости,
поэтому $FQ$ --- её след на основании.

Далее
\[
P=FQ\cap CD.
\]

Точки $P$ и $N$
принадлежат грани $SCD$,
значит прямая $PN$
лежит в этой грани
и пересекает $SD$
в искомой вершине сечения $R$:
\[
R=PN\cap SD.
\]

Если на каком-либо шаге
нужные прямые параллельны,
выбирают другой вспомогательный след
или используют параллельность
как часть построения.

% ============================================================
% БЛОК-СХЕМА ГЕОМЕТРИИ
% ============================================================

\clearpage
\thispagestyle{empty}

\begin{center}
{\Large\bfseries\textcolor{darkaccent}{
Алгоритм: построение сечения по заданным точкам
}}
\end{center}

\vspace{3mm}

\begin{center}
\begin{tikzpicture}[
  node distance=3.5mm and 13mm
]

\node[
  flowstart
]
(start)
{Даны точки секущей плоскости $\alpha$};

\node[
  flowaction,
  below=of start
]
(faces)
{
Для каждой точки отметьте рёбра
и грани,
которым она принадлежит
};

\node[
  flowdecision,
  below=5mm of faces
]
(pair)
{Есть две точки $\alpha$ на одной грани?};

\node[
  flowaction,
  below=6mm of pair
]
(join)
{
Соедините их
и получите линию пересечения
$\alpha$ с этой гранью
};

\node[
  flowaction,
  right=3mm of pair,
  text width=3.5cm
]
(trace)
{
Постройте след или проекции
и получите линию пересечения
с удобной плоскостью
};

\node[
  flowaction,
  below=5mm of join
]
(newpoint)
{
Продлите найденную линию
до ребра или его продолжения
и получите новую точку
секущей плоскости
};

\node[
  flowaction,
  below=of newpoint
]
(nextface)
{
Перейдите к соседней грани,
содержащей новую точку
};

\node[
  flowdecision,
  below=5mm of nextface
]
(closed)
{
Все стороны сечения построены
и многоугольник замкнут?
};

\node[
  flowaction,
  below=5mm of closed
]
(verify)
{
Проверьте каждую сторону:
её концы должны лежать
в одной грани
};

\node[
  flowend,
  below=of verify
]
(end)
{Сечение построено};

\draw[flowarrow]
(start) -- (faces);

\draw[flowarrow]
(faces) -- (pair);

\draw[flowarrow]
(pair)
--
node[flowlabel,right]{да}
(join);

\draw[flowarrow]
(pair.east)
--
node[flowlabel,above]{нет}
(trace.west);

\draw[flowarrow]
(join) -- (newpoint);

\draw[flowarrow]
(trace.south)
|-
(newpoint.east);

\draw[flowarrow]
(newpoint) -- (nextface);

\draw[flowarrow]
(nextface) -- (closed);

\draw[flowarrow]
(closed)
--
node[flowlabel,right]{да}
(verify);

\draw[flowarrow]
(closed.west)
--
++(-1.0,0)
|-
node[
  flowlabel,
  pos=0.16,
  left
]
{нет}
(pair.west);

\draw[flowarrow]
(verify) -- (end);

\end{tikzpicture}
\end{center}

\clearpage

% ============================================================
% ТИПИЧНЫЕ ОШИБКИ
% ============================================================

\section*{4. Типичные ошибки и самопроверка}
\sectionline

\renewcommand{\arraystretch}{1.28}

\begin{tabularx}{
  \linewidth
}{
  @{}
  >{\raggedright\arraybackslash}p{0.34\linewidth}
  X
  @{}
}

\toprule

\textbf{Риск}
&
\textbf{Контрольный вопрос}
\\

\midrule

Корень выбран по догадке
&
Проверено ли значение $P(r)$ прямой подстановкой?
\\

Пропущен коэффициент $0$
&
Выписаны ли все степени от старшей до $x^0$?
\\

В Горнере получен ненулевой остаток
&
Верен ли кандидат $r$
и нет ли арифметической ошибки
в цепочке «умножить --- сложить»?
\\

Используется слишком длинный путь
&
Проверены ли общий множитель,
группировка,
формулы
и удобная замена?
\\

Соединены точки из разных граней
&
Можно ли назвать одну грань,
содержащую обе точки?
\\

Пересечены скрещивающиеся прямые
&
Установлено ли,
что обе прямые принадлежат одной плоскости?
\\

Появилась вспомогательная точка вне фигуры
&
Понятно ли,
пересечением каких двух компланарных прямых
она получена
и для перехода на какую грань используется?
\\

\bottomrule

\end{tabularx}

% ============================================================
% ТРЕНИРОВОЧНЫЙ БЛОК
% ============================================================

\clearpage

\section*{5. Тренировочный блок}
\sectionline

\textbf{Задание 1.}

Решите уравнение,
используя подбор целого корня
и схему Горнера:
\[
x^3-2x^2-5x+6=0.
\]

Запишите допустимые целые кандидаты
и таблицу Горнера
для первого найденного корня.

\vspace{6mm}

\textbf{Задание 2.}

Решите уравнение
с помощью схемы Горнера:
\[
x^4+10x^3+25x^2-36=0.
\]

Перед вычислениями восстановите коэффициент
при отсутствующем члене первой степени.
После первого понижения степени
продолжите решение удобным способом.

\vspace{6mm}

\textbf{Задание 3.}

В кубе
\[
ABCDA_1B_1C_1D_1
\]
заданы точки
\[
K\in AB,
\qquad
L\in BB_1,
\qquad
M\in D_1C_1,
\]
причём
\[
AK:KB=7:13,
\qquad
BL:LB_1=13:7,
\qquad
D_1M:MC_1=2:3.
\]

Постройте сечение куба
плоскостью $KLM$.

Подпишите дополнительные вершины сечения
и для каждой стороны укажите грань,
в которой она лежит.

\begin{center}
\begin{tikzpicture}[scale=0.90]

\tkzSetUpLine[
  line width=0.85pt,
  color=linegray
]

\tkzSetUpPoint[
  color=darkaccent,
  fill=darkaccent,
  size=3
]

\tkzSetUpLabel[
  color=black,
  font=\small
]

\tkzDefPoint(0,0){A}
\tkzDefPoint(5,0){B}
\tkzDefPoint(7,1.5){C}
\tkzDefPoint(2,1.5){D}

\tkzDefPoint(0,4){A1}
\tkzDefPoint(5,4){B1}
\tkzDefPoint(7,5.5){C1}
\tkzDefPoint(2,5.5){D1}

\tkzDefPoint(1.75,0){K}
\tkzDefPoint(5,2.6){L}
\tkzDefPoint(4,5.5){M}

\tkzDrawSegments(
  A,B
  B,C
  C,D
  D,A
  A,A1
  B,B1
  C,C1
  D,D1
  A1,B1
  B1,C1
  C1,D1
  D1,A1
)

\tkzDrawPoints(K,L,M)

\tkzLabelPoint[below left](A){$A$}
\tkzLabelPoint[below](B){$B$}
\tkzLabelPoint[right](C){$C$}
\tkzLabelPoint[left](D){$D$}
\tkzLabelPoint[left](A1){$A_1$}
\tkzLabelPoint[above](B1){$B_1$}
\tkzLabelPoint[right](C1){$C_1$}
\tkzLabelPoint[above left](D1){$D_1$}
\tkzLabelPoint[below](K){$K$}
\tkzLabelPoint[right](L){$L$}
\tkzLabelPoint[above](M){$M$}

\end{tikzpicture}
\end{center}

% ============================================================
% ОТВЕТЫ
% ============================================================

\clearpage

\section*{6. Ответы}
\sectionline

\renewcommand{\arraystretch}{1.35}

\begin{tabularx}{
  \linewidth
}{
  @{}
  >{\centering\arraybackslash}p{0.08\linewidth}
  X
  @{}
}

\toprule

\textbf{№}
&
\textbf{Ответ}
\\

\midrule

1
&
$x=-2,\;1,\;3$.
Для $r=1$ схема Горнера даёт частное
$x^2-x-6$.
\\

\addlinespace

2
&
$x=-6,\;-3,\;-2,\;1$.
Строка коэффициентов начинается как
$1,\;10,\;25,\;0,\;-36$.
После $r=1$ получается
$x^3+11x^2+36x+36$.
\\

\addlinespace

3
&
Сечение --- шестиугольник
$KLQMRP$,
где
$Q\in B_1C_1$,
$R\in DD_1$,
$P\in AD$.
Порядок обхода:
$K\to L\to Q\to M\to R\to P\to K$.
Дополнительно:
$B_1Q:QC_1=7:12$,
$DR:RD_1=3:2$,
$AP:PD=7:12$.
\\

\bottomrule

\end{tabularx}

\vfill

\end{document}